\section{LA FORZA ELETTROMOTRICE NEI CIRCUITI IN MOVIMENTO}

\subsection{La legge di Faraday--Newmann per i circuiti in moto}

I fenomeni che si manifestano quando un conduttore si trova all'interno di un
campo elettromagnetico in generale sono molto complessi, in quanto il modo in
cui cariche e correnti si distribuiscono nel conduttore non \`e assegnato a
priori ma dipende, oltre che dai campi esterni, anche dai campi che esse stesse
producono: un esempio di questo tipo di problemi \`e fornito dall'effetto pelle
che, proprio per la complessit\`a del problema, abbiamo discusso solo in termini
qualitativi. Nella maggior parte dei casi che abbiamo trattato, i conduttori
erano costituiti dai fili dei circuiti percorsi da corrente per cui, almeno in
regime di correnti lentamente variabili, la distribuzione della corrente
all'interno dei fili conduttori non \`e n\'e particolarmente rilevante, n\'e
particolarmente influenzata dai campi generati dal circuito stesso.\\
A maggior ragione sono complicati i problemi riguardanti i conduttori in
movimento in un campo elettromagnetico. Tuttavia, siccome anche nella vita di
tutti i giorni abbiamo a che fare con gli effetti del moto di conduttori in
campi magnetici (motori, generatori di corrente alternata, \dots), dedichiamo
questo capitolo alla discussione degli effetti pi\`u rilevanti e delle idee
generali che stanno alla loro base.

Consideriamo un circuito conduttore $\gamma(t)$ in moto in un campo magnetico
statico $\vec B(x,y,z)$. Non \`e richiesto che il circuito si muova rigidamente:
esso non solo pu\`o cambiare posizione, ma durante il moto pu\`o anche
deformarsi. Sia $\vec v$ la velocit\`a (all'istante $t$) del generico punto di
$\gamma$; le cariche nel circuito hanno velocit\`a $\vec u+\vec v$ dove $\vec u$
\`e la velocit\`a delle cariche rispetto al circuito ed \`e non nulla solo se si
riferisce ai portatori di carica e se nel circuito scorre corrente: quindi
$\vec u$, se non \`e nulla, \`e tangente a $\gamma$ (sia $\vec u$ che $\vec v$
in generale dipendono dal punto del circuito considerato).
\fig{fig10_1}{10.1}
A causa della presenza del campo magnetico, le cariche sono soggette alla forza
di Lorentz \ $\vec F=(q/c)(\vec u+\vec v\,)\wedge\vec B$: la circuitazione di
questa forza lungo $\gamma$ (divisa per $q$) \`e per definizione la forza
elettromotrice nel circuito; siccome $\vec u\wedge\vec B$ \`e ortogonale a
$\gamma$, si ha
\begin{equation}
  \mathcal E_\gamma(t)=\oint_\gamma\frac{\vec v}{c}\wedge\vec B\cdot d\vec l.
  \tag{10.1}
\end{equation}
Moltiplichiamo ambo i membri della (10.1) per $\Delta t$ e permutiamo i fattori
nel triplo prodotto:
\begin{equation}
  (\vec v\Delta t)\wedge\vec B\cdot d\vec l
  =\vec B\cdot(d\vec l\wedge\vec v\Delta t)=\vec B\cdot d\vec a
  \tag{10.2}
\end{equation}
dove $d\vec a$ \`e l'area spazzata dal tratto $d\vec l$ del circuito
nell'intervallo di tempo $t,\,t+\Delta t$. Detta $\Delta A$ l'area spazzata
dall'intero circuito $\gamma$ nell'intervallo di tempo $t,\,t+\Delta t$ (la
superficie laterale nella figura 10.1), per la (10.1) e la (10.2) si ha
\begin{equation}
  \mathcal E_\gamma(t)\Delta t=\frac{1}{c}\Phi_{\Delta A}(\vec B\,)\,.
  \tag{10.3}
\end{equation}
Siccome \ $\dive\vec B=0$, \`e nullo il flusso di $\vec B$ uscente dal volume
racchiuso da $\Delta A$ e da due (arbitrarie) superfici che rispettivamente
hanno come contorni $\gamma(t+\Delta t)$ e $\gamma(t)$, per cui
\begin{equation}
  \Phi_{\Delta A}(\vec B\,)
  =-\Big(\Phi_{\gamma(t+\Delta t)}(\vec B\,)-\Phi_{\gamma(t)}(\vec B\,)\Big)
  \tag{10.4}
\end{equation}
e quindi dalla (10.3) e (10.4)
\begin{equation}
  \mathcal E_\gamma(t)
  =-\frac{1}{c}\frac{d\,\Phi_{\gamma(t)}(\vec B\,)}{dt}\,\cdot
  \tag{10.5}
\end{equation}
La (10.5) \`e identica alla (8.2): in ogni caso la f.e.m.\ in un circuito \`e
data dalla variazione del flusso di $\vec B$ attraverso il circuito, sia essa
dovuta al fatto che il circuito \`e fermo, ma $\vec B$ varia nel tempo (flusso
concatenato), oppure al moto del circuito in un campo magnetico statico (flusso
tagliato), o anche a tutti e due gli effetti insieme. Nelle due situazioni che
portano rispettivamente alla (8.2) e alla (10.5), le cause delle f.e.m.\ sono (o
almeno sembrano) diverse: nel primo caso la f.e.m.\ \`e dovuta all'esistenza di
un campo elettrico non conservativo, mentre nel secondo caso \`e presente solo
un campo magnetico, e la f.e.m.\ trae origine dalla forza di Lorentz sulle
cariche; si pu\`o anche dire che nel primo caso la f.e.m.\ esiste anche su
circuiti non materializzati da un conduttore, mentre nel secondo caso \`e
presente solo in un circuito conduttore.\\
Tuttavia, la distinzione che abbiamo fatto fra la f.e.m.\ in un circuito fermo e
in un circuito in moto, non \`e cos\`i essenziale come pu\`o apparire a prima
vista: supponiamo, per esempio, che il circuito si muova di moto rettilineo
uniforme (in direzione dell'asse $x$) in un campo magnetico statico
$\vec B(x,y,z)$. Nel riferimento $K'$ solidale con il circuito \`e presente un
campo elettrico $\vec E\,'$ dato dalla (6.38) (con i ruoli di $K$ e di $K'$
scambiati): se, visto che le velocit\`a dei conduttori sono sempre molto piccole
rispetto alla velocit\`a della luce, poniamo $\gamma(v)=1$, in $K'$ si ha
\begin{equation}
  \vec E\,'(x',y',z',t')=\frac{\vec v}{c}\wedge\vec B\,'(x',y',z',t')
  \approx\frac{\vec v}{c}\wedge\vec B(x'+vt',y',z')
  \tag{10.6}
\end{equation}
e
\begin{equation}
  \rot\vec E\,'=-\frac{1}{c}\frac{\partial\vec B\,'}{\partial t'}
  \approx-\frac{v}{c}\frac{\partial\vec B}{\partial x}
  \tag{10.7}
\end{equation}
quindi, in $K'$, la f.e.m.\ nel circuito \`e dovuta al campo elettrico
$\vec E\,'$ e coincide (nell'approssimazione $\gamma(v)=1$) con la
f.e.m.\ calcolata nel riferimento $K$ del laboratorio. Si noti che se
\ $\partial\vec B/\partial x=0$ \ (per esempio nel caso di un campo uniforme),
$\vec E\,'$ \`e conservativo, e non si ha variazione di flusso n\'e in $K$
(flusso tagliato), n\'e in $K'$ (flusso concatenato):
\ $\partial\vec B/\partial x\ne0$ \ \`e essenziale affinch\'e in $K'$ il campo
magnetico dipenda dal tempo e dia quindi luogo ad un campo elettrico non
conservativo.\\
Quindi, almeno nel caso di circuiti in moto uniforme, l'attribuzione della causa
della f.e.m.\ alla forza di Lorentz o ad un campo elettrico non conservativo,
non \`e un fatto intrinseco, ma \`e solo una questione di scelta del
riferimento, cio\`e di punto di vista. Alla stessa conclusione si perviene anche
quando il moto del circuito non \`e uniforme, per lo meno nel caso in cui \`e un
moto rigido.

Se $I$ \`e la corrente che circola nel circuito (in moto) $\gamma$, la
f.e.m.\ sviluppa una potenza \ $W=\mathcal E_\gamma I$, che di solito viene
(almeno in parte) dissipata per effetto Joule nella resistenza del circuito;
tuttavia \`e chiaro che la forza di Lorentz, che \`e una forza a lavoro
nullo, trasferisce complessivamente al sistema una potenza nulla, quindi il
lavoro sugli elettroni dovuto alla f.e.m.\ deve essere accompagnato da un lavoro
opposto: infatti abbiamo visto che quando nel circuito circola corrente \`e solo
una parte della forza di Lorentz \ $\vec F=(q/c)(\vec u+\vec v\,)\wedge\vec B$
\ che \`e responsabile della f.e.m., e si ha
\begin{equation}
  0=\vec F\cdot(\vec u+\vec v\,)
  \equiv\frac{q}{c}(\vec u+\vec v\,)\wedge\vec B\cdot(\vec u+\vec v\,)
  \tag{10.8}
\end{equation}
e quindi su un tratto $dl$ di circuito, se $\lambda$ \`e la densit\`a dei
portatori di carica ($I=\lambda u$),
\begin{equation}
  \frac{\lambda dl}{c}\vec v\wedge\vec B\cdot\vec u
  +\frac{\lambda dl}{c}\vec u\wedge\vec B\cdot\vec v=0\,.
  \tag{10.9}
\end{equation}
Il primo termine nella (10.9), integrato su tutto il circuito, \`e la potenza
$W_\mathrm{f.e.m.}$ sviluppata dalla f.e.m., ed il secondo (integrato sul
circuito), uguale e opposto al primo, \`e il lavoro negativo fatto (nell'unit\`a
di tempo) dalla forza di Lorentz contro le forze che mantengono in moto il
circuito o, in assenza di forze meccaniche esterne, coincide con la variazione
(negativa) dell'energia cinetica del circuito: detta $W_\mathrm{mecc}$ la
potenza meccanica della forza di Lorentz sulla spira, la (10.9) pu\`o quindi
essere riscritta come
\begin{equation}
  W_\mathrm{f.e.m.}+W_\mathrm{mecc}=0\,.
  \tag{10.10}
\end{equation}
Questa discussione riprende quella che abbiamo fatto alla fine del capitolo 7.6
per spiegare come la forza di Lorentz pu\`o mettere in moto una spira percorsa
da corrente; in effetti, la (10.10) pu\`o essere letta in due modi: se nel
circuito un generatore fa circolare una corrente, il circuito viene messo in
moto dalla forza di Lorentz a spese del generatore
\ $(W_\mathrm{mecc}>0,\ W_\mathrm{f.e.m.}<0)$; se invece il circuito viene
mantenuto in movimento da forze esterne in esso viene indotta una
f.e.m.\ $(W_\mathrm{f.e.m.}>0,\ W_\mathrm{mecc}<0)$, e possiamo concludere che
dal punto di vista energetico sono le forze meccaniche esterne che determinano
la potenza sviluppata dalla forza elettromotrice.\\
Quindi, nel primo caso si realizza un motore elettrico, mentre nel secondo si
realizza un generatore di f.e.m.

Da quanto sopra segue che possiamo dare una espressione per la forza risultante
che un campo magnetico esercita su una spira percorsa da una corrente $I$,
valida anche nel caso in cui la spira non \`e piccola.\\
Immaginiamo che la spira, a partire dalla posizione in cui vogliamo calcolare la
forza su di essa, effettui un moto rigido e traslatorio: questo deve essere
pensato come un moto virtuale, cio\`e indipendente da quello che \`e il moto
reale della spira. In corrispondenza a questo moto ``virtuale'' dalla (10.10) si
ha
\begin{equation}
  W_\mathrm{mecc}=-W_\mathrm{f.e.m.}
  =-I(t)\mathcal E_\gamma(t)=\frac{I(t)}{c}\frac{d\,\Phi(\vec B\,)}{dt}\,\cdot
  \tag{10.11}
\end{equation}
Siccome il moto virtuale della spira \`e rigido e traslatorio, la sua posizione
all'istante $t$ \`e univocamente determinata dalla posizione all'istante
iniziale e dalla legge oraria $\vec r(t)$ di un suo punto, quindi
\ $\Phi(\vec B\,)$ \ \`e funzione (solo) di $\vec r(t)$ e si ha
\begin{equation}
  \frac{d\,\Phi(\vec B\,)}{dt}
  =\frac{\partial\,\Phi(\vec B\,)}{\partial x_i}\,\frac{dx_i}{dt}
  =\grad\Phi(\vec B\,)\cdot\vec v
  \tag{10.12}
\end{equation}
e siccome tutti i punti della spira hanno la stessa velocit\`a $\vec v$,
\ $W_\mathrm{mecc}=\vec F\cdot\vec v$. Dalla (10.11), essendo $\vec v$
arbitraria (il moto virtuale \`e arbitrario), si ha
\begin{equation}
  \vec F=\frac{I}{c}\grad\Phi(\vec B\,)
  \tag{10.13}
\end{equation}
che si riduce alla (7.62) se la componente di $\vec B$ ortogonale alla spira
pu\`o essere considerata costante sull'area della spira. Si noti che la (10.13)
fornisce la forza \underline{risultante} agente sulla spira, per cui in pratica
non \`e di grande utilit\`a se la spira (o il circuito) \`e deformabile (come il
circuito che discuteremo nel prossimo esempio), dato che in questo caso ci\`o
che interessa \`e la forza sulle varie parti che costituiscono il circuito.\\
Analogamente alla (10.13), il momento risultante agente sulla spira \`e
determinato dalla derivata di $\Phi(\vec B\,)$ rispetto all'orientamento della
spira.

\paragraph{Esempi}
\begin{enumerate}
  \item Una sbarretta metallica si muove con velocit\`a uniforme $\vec v$ su due
    guide metalliche parallele, a distanza $d$, collegate fra di loro da un filo
    conduttore, come in figura. Il dispositivo \`e soggetto ad un campo
    magnetico uniforme ortogonale al piano delle guide. Sia $R$ la resistenza
    del circuito, che supponiamo tutta concentrata nella sbarretta. Vogliamo
    calcolare la f.e.m.\ nel circuito, la potenza dissipata per effetto Joule, e
    la potenza delle forze meccaniche che occorre esercitare sulla sbarretta per
    mantenerla in moto con velocit\`a $\vec v$.
    \fig{fig10_2}{10.2}
    La variazione di flusso nel circuito \`e dovuta esclusivamente (trascuriamo
    l'auto\-in\-du\-zio\-ne del circuito) al fatto che aumenta la sua area
    (deformazione del circuito). Si ha (il verso di percorrenza nel circuito \`e
    quello antiorario)
    \begin{equation}
      \mathcal E=-\frac{Bvd}{c}\,;\quad\ I=\frac{\mathcal E}{R}
      =-\frac{Bvd}{Rc}\,;\quad\ W_\mathrm{J}=\frac{B^2v^2d^2}{Rc^2}\,\cdot
      \tag{10.14}
    \end{equation}
    Per la (7.53) il campo magnetico esercita sulla sbarretta la forza
    \begin{equation}
      \vec F=\frac{I}{c}\vec d\wedge\vec B
      \tag{10.15}
    \end{equation}
    dove $\vec d$ \`e il vettore di modulo $d$ parallelo alla sbarretta, il cui
    verso coincide con quello di percorrenza del circuito (dal basso verso
    l'alto). Si noti che la forza risultante sul circuito \`e nulla, dato che le
    forze che si esercitano sui lati opposti sono opposte (e in effetti per una
    traslazione rigida del circuito il flusso di $\vec B$ non cambia), ma quella
    che a noi interessa \`e la forza sulla sbarretta mobile.\\
    $\vec F$, data dalla (10.15), \`e opposta a $\vec v$; per mantenere la
    sbarretta in moto uniforme occorre esercitare dall'esterno una forza
    $-\vec F$, la cui potenza \`e
    \begin{equation}
      W_\mathrm{est}=-\vec F\cdot\vec v=-\frac{I}{c}Bvd=\frac{B^2v^2d^2}{Rc^2}
      \tag{10.16}
    \end{equation}
    che eguaglia la potenza dissipata per effetto Joule. Il risultato \`e
    indipendente dal verso di $\vec B$ e da quello di $\vec v$
    \ ($W_\mathrm{est}$ \`e quadratico in entrambi), \underline{ma non} dal
    segno $-$ nella (10.5) \big(o nella (8.2)\big); con il segno $+$ nella
    (10.5), $W_\mathrm{est}$ sarebbe negativa: un modo per ottenere dal circuito
    sia calore Joule, sia lavoro positivo sull'esterno!

    Supponiamo ora che un generatore $G$ faccia circolare nel circuito una
    corrente $I$, mantenendola costante (``generatore di corrente''). Il campo
    magnetico esercita sulla sbarretta una forza data dalla (10.15): preso
    l'asse $x$ orientato da sinistra verso destra, e supposto che per $t=0$ la
    sbarretta sia ferma si ha
    \begin{equation}
      F_x=\frac{IBd}{c}\,;\qquad v_x=\frac{F_x}{m}t=\frac{IBd}{mc}t
      \tag{10.17}
    \end{equation}
    cio\`e la sbarretta viene messa in moto e quindi il sistema \`e
    sostanzialmente un (rudimentale) motore elettrico. Da dove proviene
    l'energia cinetica della sbarretta? Dalla discussione fatta sopra risulta
    che l'energia deve essere fornita dal generatore $G$: infatti la
    f.e.m.\ indotta nel circuito \`e data dalla (10.14) e il generatore $G$, per
    mantenere $I$ costante, deve aumentare di $-\mathcal E$ la propria f.e.m.,
    quindi fornire (in pi\`u) una potenza
    \begin{equation}
      W=-\mathcal EI=\frac{I}{c}Bvd
      =F_xv_x=\frac{d}{dt}\big(\frac{1}{2}mv^2\big).
      \tag{10.18}
    \end{equation}
  \item Una spira di area $\mathcal A$ ruota con velocit\`a angolare costante
    $\omega$ attorno ad un asse ortogonale ad un campo magnetico uniforme. Se
    \ $\theta(t)=\omega t$ \ \`e l'angolo fra la normale al piano della spira e
    $\vec B$, si ha
    \fig{fig10_3}{10.3}
    \[
      \Phi(t)=\mathcal AB\cos\omega t\,;\quad\ \Rightarrow
      \quad\ \mathcal E(t)=\frac{\mathcal AB\omega}{c}\sen\omega t\,.
    \]
    In questo caso quindi nella spira abbiamo una f.e.m.\ alternata.\\
    Se ad esempio prendiamo \ $\mathcal A=400\,\mathrm{cm}^2$, \ $B=10^3$\,G,
    \ $\omega=2\pi\cdot50\,\mathrm{s}^{-1}$ (c.a.\ a 50\,Hz), si ottiene
    \begin{equation}
      \mathcal E(t)=\mathcal E_0\sen\omega t\,,
      \qquad\mathcal E_0=1.26\,\mathrm{V}\,.
      \tag{10.19}
    \end{equation}
    La f.e.m.\ pu\`o essere aumentata di un fattore $N$ prendendo, anzich\'e una
    sola spira, un avvolgimento di $N$ spire: su questi principi si basa il
    funzionamento di un alternatore (generatore di f.e.m.\ alternata).\\
    Sia ora $R$ la resistenza della spira (oppure la resistenza del circuito che
    utilizza la f.e.m.\ generata nella spira); la potenza dissipata per effetto
    Joule \`e
    \begin{equation}
      W_\mathrm{J}=\frac{\mathcal E^2}{R}
      =\frac{\mathcal A^2B^2\omega^2}{Rc^2}\sen^2\omega t\,.
      \tag{10.20}
    \end{equation}
    Il campo magnetico esercita sulla spira un momento $\vec M$ parallelo
    all'asse di rotazione (asse $z$), che per le (7.49) e (7.59) \`e dato da
    \begin{equation}
      M_z=\frac{I}{c}(\vec{\mathcal A}\wedge\vec B)_z
      =-\frac{I}{c}\mathcal AB\sen\omega t\,.
      \tag{10.21}
    \end{equation}
    e quindi il motore che mantiene in rotazione la spira deve fornire un
    momento $-\vec M$ e una potenza
    \begin{equation}
      W_\mathrm{motore}=-\vec M\cdot\vec\omega
      =-M_z\omega=\frac{\mathcal A^2B^2\omega^2}{Rc^2}\sen^2\omega t
      \tag{10.22}
    \end{equation}
    uguale a quella dissipata per effetto Joule.
\end{enumerate}

\subsection{Conduttori in moto in un campo magnetico}

Consideriamo un conduttore in moto in un campo magnetico statico: per esempio
una sbarra conduttrice in moto traslatorio, oppure un disco di rame che ruota
attorno al proprio asse; i circuiti indeformabili in moto considerati nel
paragrafo precedente sono un caso particolare. Supporremo che il conduttore sia
di un materiale, come ad esempio il rame o l'alluminio, ma non un materiale
ferromagnetico, che non altera apprezzabilmente il campo magnetico.\\
Sugli elettroni di conduzione agisce la forza di Lorentz, per cui pu\`o
verificarsi l'una o l'altra delle due seguenti situazioni: le cariche libere nel
conduttore si ridistribuiscono in modo da generare un campo elettrico che in
ogni punto interno al conduttore annulla esattamente la forza di Lorentz (per
unit\`a di carica), e in questo caso all'interno (e sulla superficie) del
conduttore le cariche sono in equilibrio. Questa situazione \`e analoga a quella
di un conduttore posto in un campo elettrostatico. Se invece non \`e possibile
una ridistribuzione delle cariche che garantisca all'interno del conduttore
l'annullamento della forza totale agente sulle cariche libere, nel conduttore si
ha un continuo moto di cariche: ``correnti parassite'' o ``correnti di
Foucault''.\\
Ovviamente qui non ci occupiamo del moto transitorio delle cariche per portarsi
nella configurazione di equilibrio (se esiste), o per instaurare le correnti
parassite permanenti, ma della situazione di regime.\\
Per capire quale delle due citate situazioni si verifica, basta considerare che
un moto (permanente) di cariche richiede l'esistenza di f.e.m.\ indotte: se,
considerato un qualsiasi circuito chiuso all'interno del conduttore (e con esso
solidale), \`e nulla la f.e.m.\ data dalla (10.5), non sono possibili correnti
indotte nel conduttore, e le cariche sono in quiete rispetto al conduttore;
l'esistenza, invece, di circuiti nei quali la (10.5) d\`a una f.e.m.\ non nulla
comporta l'esistenza delle correnti di Foucault.\\
Anzich\'e affrontare il problema in generale, esaminiamo alcuni casi
particolarmente semplici ma significativi.
\begin{enumerate}
  \item Un conduttore si muove di moto traslatorio con velocit\`a costante
    $\vec v$ in un campo magnetico uniforme $\vec B$.\\
    Il problema, come quello del circuito in moto uniforme discusso nel
    paragrafo precedente, pu\`o essere affrontato sia nel riferimento $K$ del
    laboratorio, sia nel riferimento $K'$ solidale al conduttore.\\
    Nel laboratorio: siccome $\vec B$ \`e uniforme, \`e nulla la variazione di
    flusso attraverso qualsiasi circuito solidale al conduttore, quindi non si
    hanno correnti di Foucault. Le cariche di conduzione si ridistribuiscono in
    modo da creare \underline{all'interno} del conduttore un campo elettrico
    \begin{equation}
      \vec E=-\frac{\vec v}{c}\wedge\vec B.
      \tag{10.23}
    \end{equation}
    Siccome $\vec E$ dato dalla (10.23) \`e uniforme (dentro al conduttore), si
    ha
    \begin{equation}
      \dive\vec E=0\quad\ \Rightarrow\quad\ \rho=0
      \tag{10.24}
    \end{equation}
    e quindi le cariche si distribuiscono sulla superficie: in generale non \`e
    possibile determinare la distribuzione superficiale $\sigma$. All'esterno
    del conduttore si ha il
    campo (non uniforme) generato da $\sigma$. Se $\vec v$ \`e parallelo a
    $\vec B$, non succede nulla: \ $\vec E=0$, \ $\sigma=0$.

    In $K'$: per le (6.35) \big(o la (6.38)\big), in assenza del conduttore, \`e
    presente un campo elettrico dato da \big($\vec v=(v,0,0)$\big)
    \begin{equation}
      \vec E\,'=(0,\,-\beta\gamma B_z\,,\,\beta\gamma B_y)
      =\gamma(v)\frac{\vec v}{c}\wedge\vec B.
      \tag{10.25}
    \end{equation}
    $\vec E\,'$ \`e un campo statico conservativo (\`e addirittura uniforme),
    quindi questo \`e ora un problema di elettrostatica: le cariche si
    distribuiscono sulla superficie con distribuzione $\sigma'$ in modo da
    generare all'interno del conduttore un campo opposto ad $\vec E\,'$;
    $\sigma$ e $\sigma'$ coincidono se poniamo $\gamma(v)=1$ (in questa
    approssimazione anche la forma del conduttore \`e la stessa in $K$ e in
    $K'$).\\
    Nel laboratorio all'interno del conduttore \`e presente sia il campo
    elettrico che il campo magnetico; in $K'$, invece, dentro al conduttore \`e
    presente solo il campo magnetico: cambia la descrizione, ma la fisica \`e la
    stessa, cio\`e la forza sugli elettroni \`e nulla indipendentemente dal
    riferimento.
\end{enumerate}

\paragraph{Esempi}
\begin{enumerate}
  \item Il conduttore \`e una lastra ``infinita'' (cio\`e trascuriamo gli
    effetti di bordo) con un certo spessore, parallela al piano $z=0$, il campo
    magnetico $\vec B$ \`e parallelo all'asse $y$ e $\vec v$ all'asse $x$
    (figura 10.4).
    \fig{fig10_4}{10.4}
    In $K'$ (riferimento solidale alla lastra) dalla (10.25)
    \begin{equation}
      \vec E\,'=(0,\,0,\,\beta\gamma B)
      \tag{10.26}
    \end{equation}
    quindi sulle superfici della lastra si hanno le distribuzioni $\sigma'$ e
    $-\sigma'$, con
    \begin{equation}
      \sigma'=\frac{E\,'}{4\pi}=\frac{\beta\gamma B}{4\pi}
      \tag{10.27}
    \end{equation}
    che generano (solo) all'interno della lastra un campo opposto ad
    $\vec E\,'$.\\
    Nel laboratorio, dalla (10.23) si ottiene
    \begin{equation}
      \sigma=\frac{E}{4\pi}=\frac{\beta B}{4\pi}
      \tag{10.28}
    \end{equation}
    mentre dalla (10.27) e dalla (6.3) (in $K'$ le cariche sono in quiete) si
    ottiene
    \begin{equation}
      \sigma=\gamma\sigma'=\frac{\beta\gamma^2B}{4\pi}
      \tag{10.29}
    \end{equation}
    che differisce per il fattore $\gamma^2$ dalla (10.28). Il motivo della
    discrepanza \`e dovuto al fatto che in $K$ le cariche indotte sulla
    superficie del conduttore sono in moto, quindi generano a loro volta un
    campo magnetico (campo magnetico autoindotto) del quale non abbiamo tenuto
    conto nello scrivere la (10.23), e quindi la (10.28); quindi il risultato
    ottenuto in $K'$ \`e esatto, mentre quello in $K$ \`e solo approssimato. In
    questo caso \`e possibile risolvere esattamente il problema anche in $K$, ma
    in generale (conduttore di forma arbitraria) \`e certamente pi\`u semplice
    (in linea di principio) risolvere il problema nel riferimento solidale al
    conduttore. Tuttavia, il moto in $K$ delle cariche indotte sulla superficie
    del conduttore genera un campo magnetico di ordine $\beta\sigma$, e quindi
    di ordine $\beta^2B$: siccome il moto dei conduttori \`e lento (rispetto a
    $c$), in pratica \`e pi\`u che lecito trascurare termini di ordine
    $(v/c)^2$, ed in questa approssimazione le trattazioni in $K$ e in $K'$ sono
    equivalenti.

  \item Un disco conduttore ruota con velocit\`a angolare $\omega$ attorno al
    proprio asse, parallelo ad un campo magnetico uniforme $\vec B$.\\
    Anche in questo caso il flusso di $\vec B$ tagliato da qualsiasi circuito
    solidale al conduttore non varia nel tempo, quindi non si hanno
    f.e.m.\ indotte. Se trascuriamo il campo autoindotto (che \`e di ordine
    $\beta^2B$), nel riferimento del laboratorio all'interno del disco si ha
    \ $(\vec v=\vec\omega\wedge\vec r\,)$
    \begin{equation}
      \vec E=-\frac{\vec v}{c}\wedge\vec B=-\frac{\omega B}{c}\vec\varrho
      \tag{10.29}
    \end{equation}
    dove $\vec\varrho$ \`e il vettore nel piano ortogonale a $\vec B$ dall'asse
    di rotazione al generico punto del conduttore. Quindi nei punti interni al
    disco si ha una distribuzione di carica $\rho$ data da
    \begin{equation}
      \rho=\frac{1}{4\pi}\dive\vec E=-\frac{\omega B}{2\pi c}\,\cdot
      \tag{10.30}
    \end{equation}
    Poich\'e il conduttore \`e neutro, deve essere presente anche una
    distribuzione superficiale $\sigma$, che per\`o non pu\`o essere calcolata
    esplicitamente (se il conduttore fosse un cilindro infinito, $\sigma$
    sarebbe distribuita in modo uniforme sulla superficie laterale, e quindi
    facilmente calcolabile).\\
    Il campo elettrico, che siamo in grado di calcolare solo all'interno del
    disco, \`e (esattamente) conservativo: infatti sia il campo magnetico
    esterno $\vec B$, sia quello (non noto) generato dal moto della
    distribuzione di carica (il campo autoindotto), sono indipendenti dal tempo
    (la corrente dovuta al moto di $\rho$ e $\sigma$ \`e stazionaria grazie alle
    simmetrie del sistema), per cui \ $\rot\vec E=0$.\\
    In generale, se il corpo rotante non \`e un disco possiamo solo garantire
    che il campo elettrico dovuto alla ridistribuzione delle cariche \`e
    conservativo (ma non statico) solo se trascuriamo il campo magnetico
    autoindotto che, a causa del moto delle cariche indotte, potrebbe dipendere
    dal tempo.

    Il disco rotante \`e simile ad una pila a circuito aperto: come nella pila,
    dentro al disco la forza di Lorentz \`e bilanciata dalla presenza di un
    campo elettrico, ma se il circuito viene chiuso esternamente al disco si ha
    un passaggio di corrente: infatti, supponiamo che il disco rotante sia
    collegato ad una resistenza $R$ per mezzo di due conduttori, uno fissato
    sull'asse del disco, e l'altro a contatto strisciante con il bordo (figura
    10.5); il circuito chiuso all'interno del disco da un qualsiasi cammino
    solidale al disco fra il centro ed il contatto strisciante (per esempio
    $A-C-B$), non \`e pi\`u un circuito interamente solidale al disco nel quale
    quindi la circuitazione della forza di Lorentz \`e nulla (o,
    equivalentemente, nel quale il flusso tagliato non cambia), ma \`e un
    circuito deformabile e in esso si ha una f.e.m.\ data da
    \fig{fig10_5}{10.5}
    \begin{equation}
      \mathcal E=\frac{1}{c}\int_A^B\vec v\wedge\vec B\cdot d\vec l
      =\frac{\omega Br_0^2}{2c}\equiv-\frac{1}{c}\frac{d\Phi(\vec B\,)}{dt}
      \tag{10.31}
    \end{equation}
    dove $r_0$ \`e il raggio del disco.\\
    Ora nel circuito si ha dissipazione di energia per effetto Joule, e quindi
    il motorino che mantiene in rotazione il disco deve erogare una potenza
    corrispondente: infatti, se $R$ \`e la resistenza del circuito, sulla
    corrente \ $I=\mathcal E/R$ \ che circola nel disco il campo magnetico
    esercita una forza il cui momento assiale \`e
    \begin{equation}
      M_B=-\int_0^{r_0}\frac{IB}{c}\,r\,dr
      =-\frac{IBr_0^2}{2c}=-\frac{\omega B^2r_0^4}{4Rc^2}
      \tag{10.32}
    \end{equation}
    quindi il motorino deve esercitare un momento \ $M=-M_B$ \ ed una potenza
    \begin{equation}
      W=-M_B\omega=\frac{\omega^2B^2r_0^4}{4Rc^2}
      =\frac{\mathcal E^2}{R}=W_\mathrm{J}.
      \tag{10.33}
    \end{equation}
    Anche in questo caso il motorino \`e energeticamente responsabile della
    f.e.m.\ nel circuito; viceversa, se nel circuito un generatore $G$ fa
    circolare corrente, il momento esercitato dal campo magnetico sulla corrente
    mette in rotazione il disco che in tal modo funziona da motore, anzich\'e da
    generatore di f.e.m..

  \item Un cilindro pieno conduttore ruota con velocit\`a angolare $\omega$
    attorno al proprio asse, ortogonale ad un campo magnetico uniforme $\vec B$.
    \fig{fig10_6}{10.6}
    Il flusso tagliato da un qualsiasi circuito solidale al disco e non giacente
    in un piano ortogonale all'asse di rotazione varia nel tempo, quindi nel
    cilindro si hanno correnti indotte: infatti, una eventuale ridistribuzione
    delle cariche di conduzione pu\`o solo generare un campo elettrico
    conservativo $(\partial\vec B/\partial t=0)$, che non contribuisce alla
    f.e.m., determinata quindi esclusivamente dalla forza di Lorentz.
    Contrariamente ai casi precedenti, ora non siamo in grado di calcolare il
    campo elettrico dentro al conduttore, in quanto viene meno la condizione che
    le cariche siano in quiete. Di conseguenza, non possiamo neppure calcolare
    la distribuzione della corrente $\vec J(x,y,z)$ all'interno del cilindro, in
    quanto questa \`e determinata anche dal campo elettrico (conservativo)
    generato dalla ridistribuzione delle cariche: in questo caso, in cui la
    forza sugli elettroni \`e dovuta sia al campo elettrico che a quello
    magnetico, al posto della (3.17) abbiamo
    \begin{equation}
      \vec J(x,y,z)=\sigma_c\big(\vec E+\frac{\vec v}{c}\wedge\vec B\,\big)
      \tag{10.34}
    \end{equation}
    (si noti che $E$ \`e di ordine $\beta B$, e quindi esercita sugli elettroni
    una forza dello stesso ordine di quella dovuta al campo esterno $\vec B$,
    mentre il campo magnetico autoindotto \`e di ordine $\beta^2B$, e perci\`o
    pu\`o essere trascurato). Possiamo farci un'idea qualitativa della
    distribuzione di $\vec J$ osservando che
    \begin{equation}
      \vec v\wedge\vec B=(\vec\omega\wedge\vec\varrho\,)\wedge\vec B
      =(\vec\omega\cdot\vec B\,)\,\vec\varrho
      -(\vec B\cdot\vec\varrho\,)\,\vec\omega=-xB\,\vec\omega
      \tag{10.35}
    \end{equation}
    quindi il termine magnetico nella (10.34) produce correnti parallele
    all'asse di rotazione, in entrambi i versi: il campo elettrico consente alle
    linee di forza di $\vec J$
    di richiudersi. Grosso modo, \`e come se il cilindro fosse costituito da
    tante spire rettangolari nei piani passanti per l'asse di rotazione.\\
    Siccome \ $\dive\vec J=0$, dalla (10.34) e (10.35) si ha che dentro il
    cilindro \ $\dive\vec E=0$, e quindi le cariche che generano $\vec E$ sono
    distribuite sulla superficie.
\end{enumerate}

Se il motorino che mantiene in rotazione il cilindro esercitando su di esso un
momento viene scollegato, il cilindro finisce per arrestarsi; se guardiamo le
cose dal riferimento solidale al cilindro, abbiamo un \textbf{campo magnetico
rotante}:
\begin{equation}
  B_x=B\cos\omega t\,,\quad\ B_y=B\sen\omega t\,,\quad\ B_z=0
  \tag{10.36}
\end{equation}
e il motorino (solidale al riferimento del laboratorio) deve esercitare un
momento per mantenere fermo il cilindro, che altrimenti (se il motorino non ci
fosse) verrebbe trascinato dal campo.\\
Quindi, tornando nel riferimento del laboratorio, se generiamo un campo
magnetico rotante, il cilindro deve mettersi in rotazione per seguire il campo,
ed in questo modo si realizza un motore (\textbf{\textit{motore ad
induzione}}).\\
Il campo magnetico rotante si ottiene sovrapponendo i campi $B_x(t)$ e $B_y(t)$
dati dalle (10.36), cio\`e due campi ortogonali che oscillano sfasati di
90$^\circ$ (in pratica si utilizzano tre campi disposti a 120$^\circ$, sfasati
di 120$^\circ$): il campo variabile induce \big(questa volta grazie alla
(8.2)\big) forze elettromotrici, e quindi correnti parassite, dentro al
cilindro, ed il momento esercitato dal campo su queste correnti mette in
rotazione il cilindro.
